\UTIResumenPagina
  {english}
  {Abstract}
  {This degree integration project addresses the accessible update and modernization of the official website of the Faculty of Engineering at Universidad Tecnológica Indoamérica (\url{https://facultad.indoamericatech.com/}), focused on ensuring autonomous digital access for people with visual disabilities (blindness and low vision). Based on an empirical technical baseline audit conducted with Google Lighthouse v13.6.0 that revealed color contrast deficiencies, heading hierarchy skips, and accessible name mismatches, the human-centred design methodology (ISO 9241-210:2019) was applied in alignment with WCAG 2.2 Level AA guidelines. Three priority modules were designed and implemented using a modern architecture based on Astro 7 and React 19: Academic Programs and Curricula, Student Procedures and FAQ, and an AI-powered Accessible Conversational Assistant integrated via keyboard shortcut (\texttt{Alt+A}) and dynamic announcements for the NVDA screen reader through \texttt{aria-live="polite"}, complemented by an accessibility toolbar. Empirical validation proved the technical resolution of identified barriers and a significant enhancement in autonomous navigation, achieving high user satisfaction and establishing a fully inclusive higher education digital platform.}
  {Keywords}
  {Artificial intelligence, Human-centred design, Visual disability, WCAG 2.2, Web accessibility.}
